\documentclass[11pt]{article}

\usepackage{amsmath,amssymb,amsthm}
\usepackage[ruled,vlined,linesnumbered]{algorithm2e}
\usepackage{hyperref}
\usepackage{booktabs}
\usepackage{natbib}
\usepackage{fullpage}
\usepackage{microtype}
\usepackage{bm}

\newtheorem{theorem}{Theorem}[section]
\newtheorem{lemma}[theorem]{Lemma}
\newtheorem{corollary}[theorem]{Corollary}
\newtheorem{proposition}[theorem]{Proposition}
\newtheorem{definition}[theorem]{Definition}
\newtheorem{remark}[theorem]{Remark}

\DeclareMathOperator{\poly}{poly}
\DeclareMathOperator{\erfc}{erfc}
\DeclareMathOperator{\corr}{corr}
\newcommand{\bigO}{\mathcal{O}}
\newcommand{\ip}[2]{\langle #1,\, #2 \rangle}

\title{Fast Tournament Equity Computation via \\
       Generating-Function Quadrature and FFT-Accelerated Subproduct Trees}

\author{
  Sam Rosenstrauch \\
  \texttt{sarose550@gmail.com}
}

\date{\today}

\begin{document}
\maketitle

\begin{abstract}
The Independent Chip Model (ICM) maps poker tournament chip stacks to
expected prize money, but exact computation requires summing over all
$n!$ elimination orderings. I present a deterministic algorithm that
reformulates ICM equity as a one-dimensional integral over
generating-function coefficients, evaluated by Gaussian quadrature
($Q = 256$ nodes, relative error ${<}\,5 \times 10^{-12}$). The
central challenge---computing leave-one-out polynomial products for all
$n$ players---is solved by an FFT-accelerated subproduct tree whose
propagation phase is the adjoint of its build phase, reducing cost from
$\bigO(nk)$ to $\bigO(n \log^2 k)$ per quadrature point. A roofline
cost model with harmonic-mean bandwidth blending across cache levels
dispatches automatically between a batched linear engine and a hybrid
block-tree engine. Implementations tuned for Apple M3~Max,
AMD Zen~4 (with AOCL-FFTW), and NVIDIA B200 GPU demonstrate the
approach at scale: the GPU computes all 1,572,864 player equities with
$k = n$ in under one second.
\end{abstract}

\section{Introduction}\label{sec:intro}

Multi-table poker tournaments are structured competitions in which
players begin with equal chip stacks and are progressively eliminated as
they lose their chips. A payout structure $(\pi_1, \pi_2, \ldots,
\pi_k)$ specifies the prize awarded to the player finishing in each
position, with $k \leq n$ positions receiving nonzero prizes. The
\emph{Independent Chip Model} (ICM) provides a principled mapping from a
player's chip stack to their expected tournament equity---the expected
prize money they will receive~\citep{malmuth2001,harville1973}.

Despite its central role in tournament strategy, exact ICM computation
is fundamentally difficult because the number of possible elimination
orderings grows factorially. For $n = 20$ players there are $2.4 \times
10^{18}$ orderings to consider. Bitmask dynamic programming reduces
this to $\bigO(n \cdot 2^n)$ time by iterating over subsets of surviving
players, but the exponential state space limits practical computation to
roughly $n \leq 23$. Monte Carlo sampling provides approximate solutions
for larger fields, but error shrinks only as $\bigO(1/\sqrt{N})$, making
high-precision results prohibitively expensive.

The key observation that enables a better approach is that each player's
ICM equity can be expressed as a one-dimensional integral rather than an
exponential sum. This follows from the \emph{exponential clock framing}:
the Malmuth--Harville elimination process is equivalent to assigning each
player an independent exponential random variable and eliminating players
in order of increasing firing time. Under this framing, the combinatorial
sum collapses into an integral that depends only on the chip stacks.

This integral is evaluated by Gaussian quadrature with $Q = 256$ nodes,
achieving exponential convergence to double-precision accuracy. The
remaining challenge is evaluating the integrand efficiently at each
quadrature point, which requires computing a degree-$k$ leave-one-out
polynomial product for each of $n$ players. I address this with an
FFT-accelerated binary subproduct tree whose propagation phase is the
adjoint of the build phase---a cross-correlation that distributes the
payout vector to all leaves simultaneously.

My contributions are:
\begin{enumerate}
  \item A rigorous derivation of ICM equities as coefficient extractions
    from a generating function evaluated via numerical quadrature
    (Section~\ref{sec:algorithm}).
  \item A hybrid engine combining a subproduct tree for inter-block
    computation with direct polynomial arithmetic for intra-block
    computation (Section~\ref{sec:algorithm}).
  \item A roofline cost model with harmonic-mean bandwidth blending that
    achieves 42/44 correct engine dispatch decisions, replacing the
    prior FMA-counting model (Section~\ref{sec:dispatch}).
  \item Platform-specific implementations for ARM64 (Apple M3~Max),
    x86-64 (AMD Zen~4 with AOCL-FFTW), and CUDA (NVIDIA B200 GPU),
    including a GPU implementation that computes $n = 1{,}572{,}864$
    equities at $k = n$ in under one second (Sections~\ref{sec:experiments}
    and~\ref{sec:gpu}).
\end{enumerate}


% ======================================================================
\section{Related Work}\label{sec:related}
% ======================================================================

\paragraph{ICM and the Malmuth--Harville model.}
The probability model underlying ICM was introduced by
\citet{harville1973} for horse-racing handicapping and adapted to poker
tournaments by \citet{malmuth2001}. Under this model, the probability
that any surviving player is eliminated next is proportional to their
chip stack. The standard exact algorithm uses bitmask dynamic programming
over all $2^n$ subsets of surviving players, running in
$\bigO(n \cdot 2^n)$ time~\citep[and references therein]{malmuth2001}.
For large fields, Monte Carlo sampling is the prevailing
approach~[AUTHOR: cite a concrete MC-ICM paper or software, e.g., ICMizer
or HoldemResources documentation], with error converging as
$\bigO(1/\sqrt{N})$.

\paragraph{Polynomial arithmetic and subproduct trees.}
The subproduct tree (also called a \emph{subresultant tree} or
\emph{product tree}) is a standard tool in computer algebra for
multi-point evaluation and interpolation.
\citet{vonzurgathen2013} provide a comprehensive treatment, including the
$\bigO(n \log^2 n)$ algorithms for computing the product of $n$ linear
factors and evaluating a polynomial at $n$ points. My use of the tree is
closest to the ``accumulating remainder'' formulation, but with the
critical difference that propagation computes cross-correlations (the
adjoint of convolution) rather than polynomial remainders. This adjoint
perspective, which I prove in Theorem~\ref{thm:tree-correct}, appears to
be new in this context.

\paragraph{FFT libraries and automatic tuning.}
FFTW~\citep{frigo2005} introduced the idea of empirical auto-tuning for
FFT plans, generating machine-specific ``codelets'' at install time.
AMD's AOCL-FFTW extends this with 636 AVX-512-specific codelets (versus
one in the upstream distribution), delivering 1.2--2$\times$ speedups at
key sizes on Zen~4 architectures~\citep{aocl2024}. Intel
MKL~\citep{intel2024mkl} provides an alternative FFT implementation;
I use dual-library dispatch (FFTW vs.\ MKL, selected per-size by
offline calibration) on x86-64. The broader philosophy of
empirically-tuned numerical libraries traces to ATLAS~\citep{whaley1998}
and SPIRAL~\citep{pueschel2005}.

\paragraph{Roofline performance models.}
The roofline model of \citet{williams2009} bounds achievable performance
by the minimum of peak compute and memory bandwidth, parameterized by
arithmetic intensity. My dispatch cost model extends this idea to
multi-level cache hierarchies using harmonic-mean bandwidth blending---a
natural consequence of modeling total transfer time as the sum of
time spent at each cache level.

\paragraph{GPU FFT and CUDA Graphs.}
NVIDIA's cuFFT library supports batched FFT execution, and the cuFFTDx
extensions~\citep{nvidia2024cufftdx} enable fused device-side FFTs
that eliminate intermediate memory round-trips. CUDA
Graphs~\citep{nvidia2024cudagraphs} amortize kernel launch overhead by
capturing and replaying entire execution graphs, which is essential for
the 40+ kernel launches in a single tree traversal.


% ======================================================================
\section{Preliminaries}\label{sec:prelim}
% ======================================================================

\subsection{Tournament Model}

\begin{definition}[Tournament instance]
A \emph{tournament instance} is a tuple $(n, \bm{S}, \bm{\pi})$ where
$n \in \mathbb{N}$ is the number of players, $\bm{S} = (S_1, \ldots, S_n)
\in \mathbb{R}_{>0}^n$ is the vector of chip stacks with $S_i > 0$ for all
$i$, and $\bm{\pi} = (\pi_1, \ldots, \pi_k) \in \mathbb{R}_{\geq 0}^k$
is the payout vector with $k \leq n$ and $\pi_1 \geq \pi_2 \geq \cdots
\geq \pi_k \geq 0$.
\end{definition}

\begin{definition}[Malmuth--Harville model]\label{def:mh}
Under the Malmuth--Harville (MH) model, the probability that player $i$ is
the first to be eliminated from a set $T \subseteq \{1, \ldots, n\}$ of
surviving players is:
\begin{equation}\label{eq:mh-elim}
  \Pr[i \text{ elim.\ first} \mid T] = \frac{S_i}{\sum_{j \in T} S_j}\,.
\end{equation}
Elimination proceeds sequentially: after the first player is eliminated, the
model is applied recursively to the remaining set $T \setminus \{i\}$ with
unchanged stacks.
\end{definition}

\begin{definition}[ICM equity]
The \emph{ICM equity} of player $i$ in a tournament $(n, \bm{S},
\bm{\pi})$ is their expected payout under the MH model:
\begin{equation}\label{eq:icm-equity}
  E_i = \sum_{m=1}^{k} \pi_m \cdot \Pr[i \text{ finishes in position } m]\,.
\end{equation}
\end{definition}

\subsection{Exponential Clock Framing}

The MH model has an equivalent continuous-time formulation that is central
to my approach. Assign each player~$j$ an independent exponential random
variable $T_j \sim \mathrm{Exp}(S_j)$ representing their elimination time.
Players are eliminated in order of increasing $T_j$. The memoryless
property ensures that the conditional probability of any surviving player
being eliminated next is proportional to their stack, recovering
Equation~\ref{eq:mh-elim}~\citep{harville1973}. Under this framing, the
probability that player~$i$ finishes in position~$r$ is:
\begin{equation}\label{eq:position-prob}
  \Pr[i \text{ finishes in position } r]
  = \int_0^\infty S_i \, e^{-S_i t}
    \!\!\!\sum_{\substack{A \subseteq [n] \setminus \{i\} \\ |A| = r-1}}
    \prod_{j \in A} e^{-S_j t}
    \prod_{j \notin A,\, j \neq i} \!\bigl(1 - e^{-S_j t}\bigr)\, dt\,,
\end{equation}
where $A$ is the set of $r-1$ players whose clocks have \emph{not} fired
by time~$t$ (they survive past player~$i$, finishing in positions $1$
through $r-1$), and the remaining $n-r$ players have been eliminated
earlier.

\subsection{Generating-Function Formulation}

Equation~\ref{eq:position-prob} can be rewritten compactly using generating
functions. Substituting $v = e^{-t}$ (so $dt = -dv/v$ and
$e^{-S_j t} = v^{S_j}$), define:

\begin{definition}[Player generating function]
For $v \in (0, 1)$, let $a_j(v) = v^{S_j}$ and
$b_j(v) = 1 - v^{S_j}$. The \emph{leave-one-out generating function} for
player $i$ is:
\begin{equation}\label{eq:gen-func}
  Q_i(x; v) = \prod_{j \neq i} \bigl(a_j(v) + b_j(v) \, x\bigr)\,.
\end{equation}
The coefficient $[x^m]\, Q_i(x; v)$ is the sum over all subsets of $m$
other players of the product of their elimination probabilities $b_j(v)$
times the remaining players' survival probabilities $a_j(v)$---exactly the
inner sum in Equation~\ref{eq:position-prob} with $m = n - r$ eliminated
players.
\end{definition}

Under the change of variables $v = e^{-t}$, Equation~\ref{eq:position-prob}
becomes:
\begin{equation}\label{eq:position-gf}
  \Pr[i \text{ finishes in position } r]
  = \int_0^1 S_i \, v^{S_i - 1}\, [x^{n-r}]\, Q_i(x; v) \, dv\,.
\end{equation}
Multiplying by $\pi_r$ and summing over $r = 1, \ldots, k$:
\begin{equation}\label{eq:equity-integral}
  E_i = \int_0^1 S_i \, v^{S_i - 1}
    \sum_{r=1}^{k} \pi_r \cdot [x^{n-r}]\, Q_i(x; v) \, dv\,.
\end{equation}
I define the per-quadrature-point inner product:
\begin{equation}\label{eq:inner-product}
  I_i(v) = \sum_{r=1}^{k} \pi_r \cdot [x^{n-r}]\, Q_i(x; v)\,.
\end{equation}

The integral is approximated by the trapezoidal rule after a change of
variables $v = \Phi(y)$ (where $\Phi$ is the standard normal CDF), which
maps $(0,1)$ to $\mathbb{R}$ and yields a rapidly decaying integrand:
\begin{equation}\label{eq:quadrature}
  E_i \approx \sum_{q=1}^{Q} w_q \cdot S_i \cdot a_i(v_q) \cdot v_q^{-1}
    \cdot I_i(v_q)\,,
\end{equation}
where $(v_q, w_q)_{q=1}^Q$ are the quadrature nodes and weights.

\subsection{Polynomial Operations and the Convolution--Correlation Adjoint}
\label{sec:conv-adj}

We work throughout with \emph{truncated polynomials} in the vector space
$\mathbb{R}[x]_{<k} \cong \mathbb{R}^k$ of polynomials of degree
$< k$, equipped with the Euclidean inner product
\begin{equation}\label{eq:poly-ip}
  \ip{f}{g} = \sum_{m=0}^{k-1} f_m \, g_m\,.
\end{equation}

\begin{definition}[Truncated polynomial multiplication (convolution)]
For $f, h \in \mathbb{R}[x]_{<k}$, define the truncated convolution
$\mu_h(f) = (f * h) \bmod x^k$ by
\begin{equation}\label{eq:trunc-conv}
  \bigl(\mu_h(f)\bigr)_m
  = \sum_{\substack{i+j=m \\ 0 \le i,j < k}} f_i \, h_j\,,
  \qquad m = 0, \ldots, k-1.
\end{equation}
The map $\mu_h \colon \mathbb{R}^k \to \mathbb{R}^k$ is linear in~$f$.
\end{definition}

The central algebraic fact underlying the propagation phase is that the
adjoint of convolution is cross-correlation.

\begin{definition}[Cross-correlation]
For $g, h \in \mathbb{R}[x]_{<k}$, the cross-correlation
$\mu_h^*(g) = \corr(g, h)$ is defined by
\begin{equation}\label{eq:corr-def}
  \bigl(\corr(g, h)\bigr)_j
  = \sum_{m=0}^{k-1-j} h_m \, g_{m+j}\,,
  \qquad j = 0, \ldots, k-1.
\end{equation}
Equivalently, $\corr(g, h)$ is the polynomial whose $j$-th coefficient is
the inner product of $h$ with the $j$-shifted truncation of~$g$.
\end{definition}

\begin{lemma}[Correlation is the adjoint of convolution]
\label{lem:adj-conv}
For all $f, g, h \in \mathbb{R}[x]_{<k}$,
\begin{equation}\label{eq:adj-identity}
  \ip{\mu_h(f)}{g} = \ip{f}{\mu_h^*(g)}\,,
  \qquad\text{i.e.,}\qquad
  \ip{f * h}{g} = \ip{f}{\corr(g, h)}\,.
\end{equation}
\end{lemma}
\begin{proof}
Expanding the left-hand side using~\eqref{eq:poly-ip}
and~\eqref{eq:trunc-conv}:
\begin{align}
  \ip{f * h}{g}
  &= \sum_{m=0}^{k-1} \Bigl(\sum_{\substack{i+j=m\\0\le i,j<k}}
      f_i\, h_j\Bigr) g_m \notag\\
  &= \sum_{i=0}^{k-1} f_i \sum_{j=0}^{k-1-i} h_j\, g_{i+j}
  \label{eq:adj-swap}\\
  &= \sum_{i=0}^{k-1} f_i \cdot \bigl(\corr(g,h)\bigr)_i
  = \ip{f}{\corr(g,h)}\,. \notag
\end{align}
The exchange of summation in~\eqref{eq:adj-swap} is justified by the
finiteness of all sums, and the upper limit $k{-}1{-}i$ on the $j$-sum
arises from the truncation $i+j < k$.
\end{proof}

\begin{remark}[Relation to the Toeplitz transpose]
\label{rem:toeplitz}
The linear map $\mu_h$ has the Toeplitz matrix representation
$M_{ij} = h_{i-j}$ (with $h_\ell = 0$ for $\ell < 0$ or $\ell \ge k$).
Lemma~\ref{lem:adj-conv} asserts that the adjoint $\mu_h^*$ is the
transposed Toeplitz matrix $M^T_{ij} = h_{j-i}$, whose action on a vector
is precisely cross-correlation.
\end{remark}

% ======================================================================
\section{Algorithm Description}\label{sec:algorithm}
% ======================================================================

\subsection{Overview}

My algorithm proceeds in three phases for each quadrature point $v_q$:
\begin{enumerate}
  \item \textbf{Compute $a_j(v_q)$:} For each player $j$, compute
    $a_j = v_q^{S_j}$ and $b_j = 1 - a_j$.
  \item \textbf{Evaluate inner products:} Compute
    $I_i(v_q) = \sum_{r=1}^{k} \pi_r \cdot [x^{n-r}]\, Q_i(x; v_q)$
    for all players $i$ (Equation~\ref{eq:inner-product}).
  \item \textbf{Accumulate:} Add the weighted contribution
    $w_q \cdot S_i \cdot a_i \cdot v_q^{-1} \cdot I_i(v_q)$ to
    $E_i$.
\end{enumerate}

Phase~2 is the computational bottleneck. I provide three engines for this
phase, with automatic cost-model-driven dispatch
(Section~\ref{sec:dispatch}).

\subsection{Engine 1: Batched Linear Forward-Backward}\label{sec:linear}

The linear engine computes all $n$ inner products in $\bigO(nk)$ time using
a forward-backward decomposition.

\begin{algorithm}[H]
\caption{Linear Forward-Backward Engine}\label{alg:linear}
\DontPrintSemicolon
\SetAlgoLined
\KwIn{$n$ players, values $a_1, \ldots, a_n$, payout vector
  $\bm{\pi}$ with $k$ positions}
\KwOut{Inner products $I_1, \ldots, I_n$}
\BlankLine
\tcp{Forward pass: compute prefix products}
$g^{(0)}_m \leftarrow \pi_{n-m}$ for $m = 0, \ldots, k-1$\;
\For{$j = 1, \ldots, n$}{
  \For{$m = 0, \ldots, k-2$}{
    $g^{(j)}_m \leftarrow a_j \cdot g^{(j-1)}_m + (1-a_j) \cdot g^{(j-1)}_{m+1}$\;
  }
  $g^{(j)}_{k-1} \leftarrow a_j \cdot g^{(j-1)}_{k-1}$\;
}
\BlankLine
\tcp{Backward pass: compute suffix products and extract equities}
$R_0 \leftarrow 1$;\quad $R_m \leftarrow 0$ for $m = 1, \ldots, k-1$\;
\For{$j = n, n-1, \ldots, 1$}{
  $I_j \leftarrow \sum_{m=0}^{k-1} g^{(j-1)}_m \cdot R_m$ \tcp*{Inner product}
  \For{$m = k-1, k-2, \ldots, 1$}{
    $R_m \leftarrow a_j \cdot R_m + (1-a_j) \cdot R_{m-1}$\;
  }
  $R_0 \leftarrow a_j \cdot R_0$\;
}
\Return{$I_1, \ldots, I_n$}
\end{algorithm}

\begin{lemma}[Correctness of the fused backward pass]\label{lem:fused}
In Algorithm~\ref{alg:linear}, the inner product on line~8 and the suffix
update on lines~9--11 can be fused into a single loop because the backward
iteration over $m$ reads $R_m$ for the dot product before overwriting it:
$R_{m-1}$ is still pristine when used in the update of $R_m$.
\end{lemma}
\begin{proof}
At iteration $m$ (decreasing), the update sets $R_m \leftarrow a_j R_m +
(1-a_j) R_{m-1}$. The value $R_{m-1}$ has not yet been updated at this
iteration (it is updated at $m-1$), and $R_m$ was already read for the dot
product accumulation $I_j \mathrel{+}= g^{(j-1)}_m \cdot R_m$. Thus the
fused loop computes both quantities correctly.
\end{proof}

\paragraph{Batched quadrature.}
To exploit SIMD parallelism, I interleave $B_Q = 8$ quadrature points
in the data layout: $a_{\text{batch}}[j \cdot B_Q + q_i]$ stores the
$q_i$-th quadrature value for player~$j$, ensuring sequential access in
the inner loops. On Apple~M3~Max this maps to 4 NEON FMA ports
$\times$ 128-bit vectors; on Zen~4 it maps natively to 512-bit AVX-512
FMA instructions.

\paragraph{L2-aware checkpointing.}
When the working set $n \cdot k \cdot B_Q \cdot 8$ bytes exceeds L2,
I employ a checkpointing scheme inspired by gradient checkpointing in
automatic differentiation. The forward pass saves checkpoints every $C$
players, where
\begin{equation}\label{eq:ckpt}
  C = \left\lfloor \frac{L_2}{k \cdot B_Q \cdot 8} \right\rfloor,
\end{equation}
sized so one segment's $g$-store fits in L2. The backward pass
recomputes each segment from its checkpoint before extracting equities,
trading $\bigO(n/C)$ recomputation passes for $\bigO(C \cdot k)$ memory.

This checkpointing is critical in the parallel (OpenMP) path as well.
An earlier implementation allocated per-thread working sets of
$n \times k \times B_Q$ doubles---at $n = 281{,}000$, $k = 50$, 16
threads, this consumed 14.4\,GB and ran at DRAM bandwidth. Applying
L2-aware checkpointing to the parallel path yielded a 13.5$\times$
speedup, demonstrating that even ``embarrassingly parallel'' workloads
require memory-hierarchy-aware design.

\subsection{Engine 2: FFT-Accelerated Subproduct Tree}\label{sec:tree}

For large $k$, the $\bigO(nk)$ linear engine becomes expensive. I replace
it with a tree-based approach whose cost is $\bigO(n \log^2 k)$.

\begin{definition}[Binary subproduct tree]
Given $n$ linear factors $P_j(x) = a_j + b_j x$ for $j = 1, \ldots, n$,
the binary subproduct tree has $L = \lceil \log_2 n \rceil + 1$ levels.
Level~0 contains the $n$ leaves $P_j$. Level~$\ell$ contains
$\lceil n / 2^\ell \rceil$ nodes, each being the product of its two
children (truncated to degree $k$):
\[
  P^{(\ell)}_i = P^{(\ell-1)}_{2i} \cdot P^{(\ell-1)}_{2i+1} \bmod x^k\,.
\]
\end{definition}

\begin{algorithm}[H]
\caption{Tree Engine: Build and Propagate}\label{alg:tree}
\DontPrintSemicolon
\SetAlgoLined
\KwIn{Leaf polynomials $P^{(0)}_1, \ldots, P^{(0)}_n$,
  payout vector $\bm{\pi}$, truncation degree $k$}
\KwOut{Leaf-level $g$-vectors $g^{(0)}_1, \ldots, g^{(0)}_n$}
\BlankLine
\tcp{Build phase (bottom-up): compute subproducts}
\For{$\ell = 1, \ldots, L-2$}{
  \For{each parent node $i$ at level $\ell$ \textbf{in parallel}}{
    $P^{(\ell)}_i \leftarrow P^{(\ell-1)}_{2i} \cdot P^{(\ell-1)}_{2i+1}
      \bmod x^k$
      \tcp*{FFT or schoolbook}
  }
}
\BlankLine
\tcp{Propagation phase (top-down): distribute payout vector}
$g^{(L-1)}_m \leftarrow \pi_{n-m}$\;
\For{$\ell = L-1, L-2, \ldots, 1$}{
  \For{each parent node $i$ at level $\ell$ \textbf{in parallel}}{
    $g^{(\ell-1)}_{2i} \leftarrow \corr\bigl(g^{(\ell)}_i,\;
      P^{(\ell-1)}_{2i+1}\bigr)$
      \tcp*{Left child via right sibling}
    $g^{(\ell-1)}_{2i+1} \leftarrow \corr\bigl(g^{(\ell)}_i,\;
      P^{(\ell-1)}_{2i}\bigr)$
      \tcp*{Right child via left sibling}
  }
}
\Return{$g^{(0)}_1[0], \ldots, g^{(0)}_n[0]$}
  \tcp*{$I_j = g^{(0)}_j[0]$}
\end{algorithm}

\subsection{Correctness of the Tree Engine via the Adjoint}\label{sec:tree-correct}

The correctness of Algorithm~\ref{alg:tree} follows from the adjoint
identity of Lemma~\ref{lem:adj-conv}. The key computational task is: for
each player~$j$, compute
\begin{equation}\label{eq:target-ip}
  I_j = \ip{\tilde\pi}{Q_j}
  = \sum_{m=0}^{k-1} \tilde\pi_m \cdot [x^m]\, Q_j(x;v)\,,
\end{equation}
where $\tilde\pi_m = \pi_{n-m}$ is the reversed payout vector and
$Q_j = \prod_{i \ne j} P_i \bmod x^k$ is the leave-one-out product.

\paragraph{The build phase as a composition of convolutions.}
The build phase constructs the full product
$R = P_1 * P_2 * \cdots * P_n \bmod x^k$
via a binary tree of truncated convolutions. Reading from root to any
leaf~$j$, the full product decomposes as
\begin{equation}\label{eq:product-decomp}
  R = Q_j * P_j \bmod x^k\,,
\end{equation}
and $Q_j$ is the product of all the \emph{sibling} polynomials encountered
on the path from leaf~$j$ to the root.

\paragraph{The propagation phase as the adjoint of the build.}
The propagation phase computes $\ip{\tilde\pi}{Q_j}$ for all $j$ by
applying the adjoint of the build phase's linear map, with $\tilde\pi$
as the seed. At each internal node, the build phase applied
$(L, R) \mapsto L * R \bmod x^k$. Holding the sibling fixed and viewing
this as a linear map in the descendant of~$j$, the adjoint is
$g \mapsto \corr(g, h)$ by Lemma~\ref{lem:adj-conv}---exactly what
Algorithm~\ref{alg:tree} implements.

\begin{theorem}[Correctness of the tree engine]\label{thm:tree-correct}
After Algorithm~\ref{alg:tree} completes, $g^{(0)}_j[0] = I_j$ for all
$j = 1, \ldots, n$.
\end{theorem}
\begin{proof}
We establish the following invariant by downward induction on $\ell$.

\textbf{Invariant.} Let $\mathrm{leaves}(i,\ell)$ denote the set of leaf
indices in the subtree rooted at node $i$ at level $\ell$, and let
$C_j^{(\ell)} \subseteq [n] \setminus \mathrm{leaves}(i,\ell)$ be the set
of all leaves \emph{outside} this subtree.  After the propagation phase
processes level $\ell$, the $g$-vector at node $i$ satisfies
\begin{equation}\label{eq:prop-invariant}
  g^{(\ell)}_i = \corr\!\Bigl(\tilde\pi,\;
    \prod_{j \in C_j^{(\ell)}} P_j \bmod x^k \Bigr).
\end{equation}

\textbf{Base case} ($\ell = L{-}1$, the root). The root's subtree
contains all $n$ players, so $C_j^{(L-1)} = \emptyset$ and the product
over the empty set is~$1$.  The initialization
$g^{(L-1)}_m = \tilde\pi_m$ is $\corr(\tilde\pi, 1) = \tilde\pi$,
satisfying~\eqref{eq:prop-invariant}.

\textbf{Inductive step.} Assume the invariant holds at level $\ell$
for node $i$.  Write $C = C_j^{(\ell)}$ and let $R_C = \prod_{j \in C}
P_j \bmod x^k$.  By the inductive hypothesis, $g^{(\ell)}_i =
\corr(\tilde\pi, R_C)$.  The left child $2i$ receives
\[
  g^{(\ell-1)}_{2i}
  = \corr\bigl(g^{(\ell)}_i,\; P^{(\ell-1)}_{2i+1}\bigr)
  = \corr\!\bigl(\corr(\tilde\pi, R_C),\; P^{(\ell-1)}_{2i+1}\bigr).
\]
We claim this equals $\corr(\tilde\pi,\; R_C \cdot P^{(\ell-1)}_{2i+1}
\bmod x^k)$.  To verify, take an arbitrary $f \in \mathbb{R}[x]_{<k}$
and compute:
\begin{align*}
  \ip{f}{g^{(\ell-1)}_{2i}}
  &= \ip{f}{\corr\bigl(\corr(\tilde\pi, R_C),\; P^{(\ell-1)}_{2i+1}\bigr)}
  \\
  &= \ip{f * P^{(\ell-1)}_{2i+1}}{\corr(\tilde\pi, R_C)}
    &&\text{(Lemma~\ref{lem:adj-conv})}
  \\
  &= \ip{(f * P^{(\ell-1)}_{2i+1}) * R_C}{\tilde\pi}
    &&\text{(Lemma~\ref{lem:adj-conv})}
  \\
  &= \ip{f * (P^{(\ell-1)}_{2i+1} * R_C)}{\tilde\pi}
    &&\text{(associativity of $*$ mod $x^k$)}
  \\
  &= \ip{f}{\corr(\tilde\pi,\; P^{(\ell-1)}_{2i+1} * R_C)}
    &&\text{(Lemma~\ref{lem:adj-conv})}\,.
\end{align*}
Since this holds for all $f$, we conclude
$g^{(\ell-1)}_{2i} = \corr(\tilde\pi,\; R_C \cdot P^{(\ell-1)}_{2i+1}
\bmod x^k)$.  But $C_j^{(\ell-1)}$ for the left child $2i$ is
$C \cup \mathrm{leaves}(2i{+}1, \ell{-}1)$, and
$P^{(\ell-1)}_{2i+1} = \prod_{j \in \mathrm{leaves}(2i+1,\ell-1)} P_j
\bmod x^k$.  Thus $R_C \cdot P^{(\ell-1)}_{2i+1} =
\prod_{j \in C_j^{(\ell-1)}} P_j \bmod x^k$, confirming the invariant
for the left child.  The right child is symmetric.

\textbf{Conclusion.} At level~$0$, leaf~$j$ has
$C_j^{(0)} = [n] \setminus \{j\}$, so
$g^{(0)}_j = \corr(\tilde\pi, Q_j)$ where $Q_j = \prod_{i \ne j} P_i
\bmod x^k$.  Evaluating at index~$0$:
\[
  g^{(0)}_j[0]
  = \bigl(\corr(\tilde\pi, Q_j)\bigr)_0
  = \sum_{m=0}^{k-1} \tilde\pi_m \cdot (Q_j)_m
  = \ip{\tilde\pi}{Q_j}
  = I_j\,.\qedhere
\]
\end{proof}

\begin{remark}[Adjoint chain rule and complexity parity]
\label{rem:adjoint-chain}
The propagation phase is an instance of reverse-mode automatic
differentiation applied to the build phase~\citep{griewank2008}.
Define the linear map $\Phi_j \colon \mathbb{R}[x]_{<k} \to
\mathbb{R}[x]_{<k}$ at each build node as $\Phi_j = \mu_{P_{\mathrm{sib}(j)}}$.
The build computes $\Phi_L \circ \cdots \circ \Phi_1$, and propagation
computes $\Phi_1^* \circ \cdots \circ \Phi_L^*$ applied to
$\tilde\pi$---the adjoint of the composition, by the chain rule
$(\Phi_L \circ \cdots \circ \Phi_1)^* = \Phi_1^* \circ \cdots \circ
\Phi_L^*$.  A standard result in reverse-mode AD states that the adjoint
of a linear program has the same arithmetic
complexity~\citep[Theorem~3.1]{griewank2008}.
\end{remark}

\paragraph{Per-level FFT vs.\ schoolbook decisions.}
At each tree level, I compare estimated FFT and schoolbook costs using
offline-calibrated data. Let $d_{\text{eff}}$ be the effective polynomial
degree at level~$\ell$. The decision rule is:
\begin{equation}\label{eq:fft-decision}
  \text{use FFT} \iff
    T_{\text{FFT}}(S^*) + T_{\text{overhead}} + (m+1)^2 \cdot \tau_{\text{FMA}}
    < (d_{\text{eff}}+1)^2 \cdot \tau_{\text{FMA}}\,,
\end{equation}
where $T_{\text{FFT}}(S^*)$ is the calibrated FFT time for optimal
smooth size $S^*$, $T_{\text{overhead}}$ is per-call overhead, $m$ is the
wrap correction count, and $\tau_{\text{FMA}}$ is the measured FMA latency.
All parameters come from offline calibration at 749 smooth sizes---not
estimated from models.

\paragraph{Cyclic FFT with wrap correction.}
Rather than padding to the next power of two, I search all 7-smooth
numbers ($2^a 3^b 5^c 7^d$) for the FFT size minimizing total cost:
\begin{equation}\label{eq:wrap-cost}
  \text{cost}(S) = T_{\text{FFT}}(S) + (m+1)^2 \cdot \tau_{\text{FMA}}\,,
  \quad m = \max(L_{\text{conv}} - S, 0)\,,
\end{equation}
where $L_{\text{conv}}$ is the required convolution length and
$m$ high-degree terms wrap and are corrected by schoolbook.

\paragraph{Paired cached correlate.}
During propagation, both children at each node require correlating the same
parent $g$-vector with their respective sibling polynomials. I share the
forward FFT of $g$ across both correlates, saving one FFT per parent node.
When build and correlate use the same FFT size, I cache the FFT of each
child during the build phase and reuse it via conjugation during
propagation.

\subsection{Engine 3: Hybrid Block-Divide}\label{sec:hybrid}

The hybrid engine partitions the $n$ players into $\lceil n/B \rceil$
blocks of size $B$ and combines direct $\bigO(B^2)$ within-block
computation with the tree engine for inter-block structure.

\begin{algorithm}[H]
\caption{Hybrid Block-Divide Engine}\label{alg:hybrid}
\DontPrintSemicolon
\SetAlgoLined
\KwIn{$n$ players with values $a_1, \ldots, a_n$, block size $B$,
  payout vector $\bm{\pi}$, degree $k$}
\KwOut{Inner products $I_1, \ldots, I_n$}
\BlankLine
\tcp{Step 1: Build block products}
\For{each block $b = 0, 1, \ldots, \lceil n/B \rceil - 1$}{
  $P_b(x) \leftarrow \prod_{j \in \text{block } b} (a_j + b_j x)$
    \tcp*{Sequential, $\bigO(B^2)$}
}
\BlankLine
\tcp{Step 2: Inter-block tree (Alg.~\ref{alg:tree}) with leaves $P_b$}
Build subproduct tree from block products $P_0, \ldots, P_{\lceil n/B \rceil - 1}$\;
Propagate $\bm{\pi}$ to obtain block-level $g$-vectors $g_b$\;
\BlankLine
\tcp{Step 3: Within-block divide}
\For{each block $b$}{
  \For{each player $j$ in block $b$}{
    Compute $Q_j^{(b)} \leftarrow P_b(x) / (a_j + b_j x)$
      via stable bidirectional synthetic division\;
    $I_j \leftarrow \sum_{m=0}^{\min(k,B)-1} g_{b,m} \cdot Q_{j,m}^{(b)}$\;
  }
}
\Return{$I_1, \ldots, I_n$}
\end{algorithm}

\begin{remark}[Bidirectional division stability]
Synthetic division of $P_b(x)$ by $(a_j + b_j x)$ is numerically unstable
when $a_j < 0.5$ (forward direction amplifies errors). I use bidirectional
division: forward when $a_j > 0.5$, backward when $a_j \leq 0.5$,
guaranteeing that the division coefficient satisfies $|c| < 1$ in all cases.
This is safe because division operates on the \emph{complete} block
product (not the truncated global product), where error amplification is
bounded by $|c|^B$ and remains innocuous for $B \leq 64$.
\end{remark}

\paragraph{Cost-model block size selection.}
The block size $B$ is selected from candidates
$B \in \{8, 16, 24, 32, 48, 64\}$ by minimizing:
\begin{equation}\label{eq:hybrid-cost}
  T_{\text{hybrid}}(B) = T_{\text{block}}(B) + T_{\text{tree}}(n/B, k)\,,
\end{equation}
where $T_{\text{block}}$ accounts for the $\bigO(nB)$ block build and
divide cost, and $T_{\text{tree}}$ is estimated using actual per-level
FFT decisions from the calibration data. Typical selections: $B = 16$ on
M3~Max, $B = 32$ on Zen~4.

\subsection{Engine Dispatch: Roofline Cost Model}\label{sec:dispatch}

Engine dispatch must compare the cost of the linear engine (memory-bound,
$\bigO(nk)$) against the hybrid engine (compute-bound in the FFT phases,
$\bigO(n \log^2 k)$). An earlier version estimated linear cost as
$2nk \cdot \tau_{\text{poly}} \cdot c_{\text{ckpt}}$, where
$\tau_{\text{poly}}$ is a single memory-bound FMA constant and
$c_{\text{ckpt}} \in \{1.0, 1.3\}$ is a binary penalty for checkpointing.
This model correctly dispatched roughly 30/44 test cases: it had no
knowledge of the cache hierarchy and could not distinguish L2-resident
from DRAM-resident working sets.

The replacement is a roofline model~\citep{williams2009} that computes
cost as bytes streamed divided by effective bandwidth at the appropriate
cache level. The key insight is that when a working set spans two cache
levels, the effective bandwidth is the \emph{harmonic mean} of the two
levels' bandwidths, weighted by the fraction of data at each level.
This follows directly from additivity of transfer times:
\begin{equation}\label{eq:harmonic-bw}
  T_{\text{total}} = \frac{B_{\text{hit}}}{W_{\text{hit}}}
    + \frac{B_{\text{miss}}}{W_{\text{miss}}}
  = \frac{B_{\text{total}}}{W_{\text{eff}}}\,,
  \qquad
  W_{\text{eff}} = \frac{1}{\alpha / W_{\text{hit}}
    + (1-\alpha) / W_{\text{miss}}}\,,
\end{equation}
where $\alpha = B_{\text{hit}} / B_{\text{total}}$ is the cache hit
fraction and $W$ denotes bandwidth.

The model requires only independently measurable hardware constants---cache
sizes and streaming bandwidths at each level---and produces dispatch
decisions directly from the physics of data movement:

\begin{algorithm}[H]
\caption{Roofline Engine Dispatch}\label{alg:dispatch}
\DontPrintSemicolon
\SetAlgoLined
\KwIn{$n, k$, hardware constants ($L_2, L_3$, $W_{L2}, W_{L3}, W_{\text{DRAM}}$)}
\KwOut{Engine choice and parameters}
\BlankLine
\tcp{Checkpoint interval: fit one segment in L2}
$C \leftarrow \lfloor L_2 \,/\, (k \cdot B_Q \cdot 8) \rfloor$\;
\BlankLine
\tcp{Linear cost: roofline at correct cache level}
\uIf{$C \geq n$}{
  $\text{bytes} \leftarrow 2nk \cdot 8$ \tcp*{Forward + backward, no ckpt}
  $T_{\text{linear}} \leftarrow \text{bytes} \,/\, W_{\text{eff}}(\text{bytes} \cdot B_Q)$\;
}
\Else{
  $T_{\text{inner}} \leftarrow (2nk \cdot 8) \,/\, W_{L2}$ \tcp*{L2-resident recompute}
  $T_{\text{outer}} \leftarrow \bigl(2(n/C) k \cdot 8 + 3 n \cdot 8\bigr) \,/\,
    W_{\text{eff}}(\text{outer bytes})$\;
  $T_{\text{linear}} \leftarrow T_{\text{inner}} + T_{\text{outer}}$\;
}
\BlankLine
\tcp{Hybrid cost: block build + calibrated FFT tree}
$B^* \leftarrow \text{select\_best\_B}(n, k)$\;
$T_{\text{hybrid}} \leftarrow T_{\text{block}}(B^*) +
  T_{\text{tree}}(n/B^*, k)$\;
\BlankLine
\lIf{$T_{\text{hybrid}} < T_{\text{linear}}$}{
  \Return{Hybrid with $B = B^*$}
}
\lElse{
  \Return{Linear batched}
}
\end{algorithm}

Validation against 44 representative $(n, k)$ pairs shows 42 correct
dispatch decisions (95.5\%), with the two errors occurring at points where
the engines are within 3--9\% of each other---well within the regime where
the choice barely matters.

% ======================================================================
\section{Theoretical Analysis}\label{sec:theory}
% ======================================================================

\subsection{Time Complexity}

\begin{theorem}[Complexity of the tree engine]\label{thm:tree-complexity}
The tree engine (Algorithm~\ref{alg:tree}) computes all $n$ inner products
in $\bigO(n \log^2 k)$ arithmetic operations per quadrature point.
\end{theorem}
\begin{proof}
Let $T_b(n)$ denote the total arithmetic cost of the build phase for $n$
leaves with truncation degree $k$.  At the root, two subtrees of
$\lfloor n/2 \rfloor$ and $\lceil n/2 \rceil$ leaves are combined by a
single truncated polynomial multiplication costing
$\bigO(\min(n, k) \cdot \log \min(n, k))$ via FFT.  This gives the
recurrence
\begin{equation}\label{eq:build-recurrence}
  T_b(n) = 2\, T_b(n/2)
  + \bigO\!\bigl(\min(n, k) \cdot \log \min(n, k)\bigr),
  \qquad T_b(1) = 0.
\end{equation}

\textbf{Case 1: $n \le k$ (below saturation).}
The recurrence becomes $T_b(n) = 2\,T_b(n/2) + \bigO(n \log n)$, which
by Case~2 of the Master theorem (with $p = 1$) yields
\begin{equation}\label{eq:below-sat}
  T_b(n) = \bigO(n \log^2 n)\,.
\end{equation}

\textbf{Case 2: $n > k$ (above saturation).}
All subproducts are capped at degree $k$, so the merge cost is
$\bigO(k \log k)$ regardless of $n$.  Unrolling for $s = \log_2(n/k)$
levels above saturation:
\begin{align}
  T_b(n) &= \frac{n}{k}\, T_b(k)
    + \bigO(k \log k) \sum_{i=0}^{s-1} 2^i \notag\\
  &= \frac{n}{k} \cdot \bigO(k \log^2 k)
    + \bigO\!\bigl(\tfrac{n}{k} \cdot k \log k\bigr)
  = \bigO(n \log^2 k)\,. \label{eq:above-sat}
\end{align}

In both cases, $T_b(n) = \bigO(n \log^2 k)$.  By
Remark~\ref{rem:adjoint-chain}, propagation has the same complexity,
so the total per-quadrature-point cost is $\bigO(n \log^2 k)$.
Over $Q$ points: $\bigO(Qn \log^2 k)$.
\end{proof}

\begin{theorem}[Complexity of the hybrid engine]\label{thm:hybrid-complexity}
The hybrid engine (Algorithm~\ref{alg:hybrid}) with block size $B$
computes all $n$ inner products in
$\bigO(nB + (n/B) \log^2 k)$ operations per quadrature point.
\end{theorem}
\begin{proof}
Block build: $n/B$ blocks of $\bigO(B^2)$ each, totaling $\bigO(nB)$.
Inter-block tree over $n/B$ leaves: $\bigO((n/B) \log^2 k)$ by
Theorem~\ref{thm:tree-complexity}. Within-block divide: $\bigO(B)$
per player, $\bigO(nB)$ total.
\end{proof}

\begin{corollary}[Optimal block size]\label{cor:opt-block}
Setting $B = \Theta(\log k)$ balances the two terms, yielding
$\bigO(n \log^2 k)$---the same asymptotic complexity as the pure tree.
In practice, hardware constraints determine $B$ via the cost model.
\end{corollary}

\begin{proposition}[Complexity of the linear engine]\label{prop:linear-complexity}
The linear forward-backward engine computes all $n$ inner products in
$\bigO(nk)$ operations per quadrature point.
\end{proposition}

\begin{remark}[Dispatch optimality]
The linear engine is asymptotically optimal for $k = \bigO(1)$ (fixed
payout depth). The hybrid/tree engines are optimal for $k = \Theta(n)$.
The roofline dispatch of Algorithm~\ref{alg:dispatch} selects the faster
engine for each $(n, k)$ pair based on measured hardware characteristics.
\end{remark}

\subsection{Numerical Stability}

\begin{lemma}[Numerical stability of the hybrid divide]\label{lem:divide}
The bidirectional synthetic division in Algorithm~\ref{alg:hybrid} has
per-step error amplification factor strictly less than~$1$ for all
$a_j \in (0, 1)$.
\end{lemma}
\begin{proof}
Forward division computes $Q_m = (P_m - (1{-}a_j) Q_{m-1}) / a_j$,
with amplification factor $(1{-}a_j)/a_j < 1$ when $a_j > 1/2$.
Backward division computes $Q_m = (P_{m+1} - a_j Q_{m+1}) / (1{-}a_j)$,
with amplification $a_j/(1{-}a_j) \leq 1$ when $a_j \le 1/2$.
Choosing forward when $a_j > 1/2$ and backward otherwise ensures errors
decay exponentially along the division chain.
\end{proof}

% ======================================================================
\section{Experimental Results}\label{sec:experiments}
% ======================================================================

\subsection{Experimental Setup}

I evaluate my implementation on three platforms:
\begin{itemize}
  \item \textbf{Apple M3~Max:} 12 performance + 4 efficiency cores,
    ARM64 with NEON (2 FP64/vector), 4 FMA ports per core,
    32\,MB shared L2 cache. FFT: FFTW3 PATIENT + Apple vDSP at
    supported sizes.
  \item \textbf{AMD Ryzen 9 7950X (Zen 4):} 16 cores,
    AVX-512 (8 FP64/vector), 2 FMA units per core,
    1\,MB per-core L2, 32\,MB shared L3. FFT: AOCL-FFTW (636 AVX-512
    codelets) + Intel MKL dual dispatch.
  \item \textbf{NVIDIA B200 GPU:} [AUTHOR: fill in SM count and
    memory bandwidth from spec sheet]. CUDA 12.8, cuFFT + cuFFTDx
    fused R2C/C2R kernels.
\end{itemize}

All CPU experiments use $Q = 256$ quadrature points, uniform chip stacks,
and report the median of 5 repetitions. Correctness is verified against
closed-form $V_1$ and $V_2$ equities and cross-checked between engines.

\subsection{CPU Performance}

\begin{table}[t]
\centering
\caption{Single-threaded execution time (ms), $Q = 256$.
  [UPDATE: re-measure with final build including AOCL-FFTW and roofline dispatch.]}
\label{tab:serial}
\begin{tabular}{@{}l rrrrrr rrrrrr@{}}
\toprule
& \multicolumn{6}{c}{\textbf{Apple M3 Max}} &
  \multicolumn{6}{c}{\textbf{AMD Zen 4 (AOCL-FFTW)}} \\
\cmidrule(lr){2-7} \cmidrule(lr){8-13}
$n$ & $k{=}10$ & $k{=}50$ & $k{=}100$ & $k{=}n/4$ & $k{=}n/2$ & $k{=}n$
    & $k{=}10$ & $k{=}50$ & $k{=}100$ & $k{=}n/4$ & $k{=}n/2$ & $k{=}n$ \\
\midrule
1024  & -- & -- & -- & -- & -- & -- & -- & -- & -- & -- & -- & -- \\
2048  & -- & -- & -- & -- & -- & -- & -- & -- & -- & -- & -- & -- \\
4096  & -- & -- & -- & -- & -- & -- & -- & -- & -- & -- & -- & -- \\
8192  & -- & -- & -- & -- & -- & -- & -- & -- & -- & -- & -- & -- \\
16384 & -- & -- & -- & -- & -- & -- & -- & -- & -- & -- & -- & -- \\
32768 & -- & -- & -- & -- & -- & -- & -- & -- & -- & -- & -- & -- \\
65536 & -- & -- & -- & -- & -- & -- & -- & -- & -- & -- & -- & -- \\
\bottomrule
\end{tabular}
\end{table}

\begin{table}[t]
\centering
\caption{16-thread parallel execution time (ms), $Q = 256$.
  [UPDATE: re-measure with parallel checkpointing fix and AOCL-FFTW.]}
\label{tab:parallel}
\begin{tabular}{@{}l rrrrrr rrrrrr@{}}
\toprule
& \multicolumn{6}{c}{\textbf{Apple M3 Max (16 threads)}} &
  \multicolumn{6}{c}{\textbf{AMD Zen 4 (16 threads)}} \\
\cmidrule(lr){2-7} \cmidrule(lr){8-13}
$n$ & $k{=}10$ & $k{=}50$ & $k{=}100$ & $k{=}n/4$ & $k{=}n/2$ & $k{=}n$
    & $k{=}10$ & $k{=}50$ & $k{=}100$ & $k{=}n/4$ & $k{=}n/2$ & $k{=}n$ \\
\midrule
1024  & -- & -- & -- & -- & -- & -- & -- & -- & -- & -- & -- & -- \\
4096  & -- & -- & -- & -- & -- & -- & -- & -- & -- & -- & -- & -- \\
8192  & -- & -- & -- & -- & -- & -- & -- & -- & -- & -- & -- & -- \\
16384 & -- & -- & -- & -- & -- & -- & -- & -- & -- & -- & -- & -- \\
65536 & -- & -- & -- & -- & -- & -- & -- & -- & -- & -- & -- & -- \\
\bottomrule
\end{tabular}
\end{table}

\paragraph{Platform comparison.}
Zen~4 benefits from AVX-512's 8-wide FP64 FMA (versus NEON's 2-wide) in
the linear engine, and from AOCL-FFTW's 636 AVX-512 codelets in the
hybrid engine. AOCL-FFTW delivers 1.2--2$\times$ speedups at FFT sizes
in the 1K--32K range over the upstream FFTW distribution, which ships only
a single AVX-512 codelet. At 749 calibrated sizes, AOCL-FFTW is
the fastest library at 637 sizes, Intel MKL at 112.

\paragraph{Parallel checkpointing.}
The parallel (OpenMP) linear engine initially allocated full
$n \times k \times B_Q$ per-thread working sets, causing DRAM spills at
moderate sizes. Adding L2-aware checkpointing to the parallel path
(matching the serial engine's existing logic) yielded a 13.5$\times$
speedup at $(n, k) = (8192, 100)$---the parallel engine now runs at
L2 bandwidth (${\sim}119$\,GB/s measured) instead of DRAM
(${\sim}27$\,GB/s).

\subsection{Quadrature Accuracy}\label{sec:accuracy}

The erfc-trap quadrature rule (trapezoidal rule after the change of
variables $v = \Phi(y)$) achieves exponential convergence in $Q$ for
smooth integrands~\citep{trefethen2014}. I compare it against tanh-sinh
quadrature~\citep{mori2001}, an alternative double-exponential scheme,
across four stack distributions:

\begin{itemize}
  \item \textbf{Uniform:} $S_i = 1$ for all $i$.
  \item \textbf{Geometric:} $S_i = 2^{i-1}$ (exponentially spaced).
  \item \textbf{Skewed:} One dominant stack ($S_1 = n/2$), rest equal.
  \item \textbf{Adversarial:} Alternating large/small stacks designed to
    stress the quadrature.
\end{itemize}

[UPDATE: Insert convergence table or figure here showing relative error
vs.\ $Q$ for both quadrature rules across all four distributions.
Data from accuracy benchmarks.]

The erfc-trap rule converges to ${\sim}10^{-12}$ relative error by
$Q = 256$ uniformly across all distributions. Tanh-sinh achieves a lower
error floor on the uniform distribution (${\sim}10^{-14}$) but fails to
converge as reliably on geometric and adversarial distributions, where
the integrand's strip of analyticity is narrower. For a general-purpose
implementation, erfc-trap is the better choice: it provides uniformly
good accuracy without distribution-dependent tuning.

\subsection{Calibration and Cost Model Validation}

Table~\ref{tab:constants} summarizes the hardware constants measured via
offline profiling. All bandwidth values are measured by streaming
benchmarks, not copied from vendor specifications.

\begin{table}[t]
\centering
\caption{Platform calibration constants.}
\label{tab:constants}
\begin{tabular}{@{}lrrl@{}}
\toprule
\textbf{Constant} & \textbf{M3 Max} & \textbf{Zen 4} & \textbf{Description} \\
\midrule
$\tau_{\text{FMA}}$ (ns) & 0.25 & 0.13 & Scalar FMA latency \\
$W_{L2}$ (GB/s) & [UPDATE] & 118.9 & Measured L2 streaming BW \\
$W_{L3}$ (GB/s) & [UPDATE] & 50.3 & Measured L3 streaming BW \\
$W_{\text{DRAM}}$ (GB/s) & [UPDATE] & 26.5 & Measured DRAM streaming BW \\
$L_2$ (bytes) & 32\,MB & 1\,MB & Per-core L2 size \\
$L_3$ (bytes) & [UPDATE] & 32\,MB & Shared L3 size \\
Typical $B$ & 16 & 32 & Cost-model block size \\
FFT backend & FFTW+vDSP & AOCL+MKL & Library dispatch \\
Dispatch accuracy & [UPDATE] & 42/44 & Correct engine selections \\
\bottomrule
\end{tabular}
\end{table}

% ======================================================================
\section{GPU Implementation: NVIDIA B200}\label{sec:gpu}
% ======================================================================

The GPU implementation targets the NVIDIA B200 and uses the same hybrid
block-tree algorithm, adapted to the GPU execution model. The linear
engine is not applicable: its sequential player-by-player structure cannot
saturate GPU parallelism. Only the tree-based engines, where all nodes at
each level are independent, map naturally to GPU execution.

\subsection{Architecture Mapping}

The tree's level-by-level structure becomes a sequence of batched kernel
launches: each tree level launches one kernel that processes all
$Q \times n_{\text{nodes}}[\ell]$ independent polynomial operations.
At the bottom levels of a tree with $n = 1{,}572{,}864$ leaves and
$Q = 256$, this yields millions of independent operations---far more than
needed to saturate the GPU's streaming multiprocessors.

\subsection{Key Optimizations}

The GPU implementation went through 12 optimization stages. I describe
the ones with the largest impact; the others are implementation details
that reduced constant factors.

\paragraph{FFT-stride layout.}
The single largest optimization was storing polynomials at their
FFT-padded stride rather than at their logical degree. In the standard
layout, cuFFT's batched interface requires gathering coefficients from
logical-length arrays into FFT-length buffers before each transform and
scattering results back afterward---two full memory passes per FFT call.
The FFT-stride layout eliminates both: polynomials are stored with
FFT-length stride from the start, so cuFFT reads and writes them
directly. This pushed the frontier past a tree-level cliff at
$n \approx 1{,}440{,}000$.

\paragraph{cuFFTDx fused R2C/C2R kernels.}
NVIDIA's cuFFTDx library~\citep{nvidia2024cufftdx} allows FFT
computation entirely within device code, fusing the forward R2C
transform, pointwise multiply, and inverse C2R transform into a single
kernel launch. This eliminates two intermediate memory round-trips per
tree node. I use cuFFTDx at sizes 2048 and 4096 (where its thread-block
FFTs are competitive with cuFFT's global algorithm) and fall back to
batched cuFFT at larger sizes.

\paragraph{Arena allocator.}
The baseline implementation made 35 separate \texttt{cudaMalloc} calls to
allocate per-level polynomial and FFT buffers. Each allocation hits the
CUDA driver, serializing execution. Replacing these with a single arena
allocation (one \texttt{cudaMalloc} for the entire tree, with per-level
offsets computed at plan time) eliminated this overhead entirely.

\paragraph{CUDA Graphs.}
The full tree traversal (build + propagation) involves 40+ kernel
launches. Capturing the entire sequence as a CUDA
Graph~\citep{nvidia2024cudagraphs} and replaying it reduces total launch
overhead to ${\sim}2.5\,\mu$s (versus ${\sim}3\,\mu$s per launch
$\times$ 40+ launches). The graph is built once per $(n, k)$
configuration and reused across all $Q$ quadrature points.

\paragraph{Q-batch pipelining.}
Rather than processing all $Q = 256$ quadrature points in a single
monolithic batch (which would require $256 \times$ more memory), I
process $Q_{\text{batch}} = 4$--$8$ points at a time, overlapping the
block-build phase of one batch with the tree propagation of the previous
batch using multi-stream execution.

\paragraph{Rejected optimizations.}
Two optimizations were tested and rejected after benchmarking:
\begin{itemize}
  \item \textbf{cuFFT callbacks:} Fusing the pointwise multiply into
    cuFFT's load/store callbacks was 2.3$\times$ \emph{slower} than the
    separate-kernel approach. The per-element callback overhead breaks
    memory coalescing, and NVIDIA has deprecated this API in favor of
    cuFFTDx.
  \item \textbf{$B = 1$ balanced tree:} Using single-player leaves
    (no block structure) was 20--56\% slower because the tree performs
    70\% more FMA operations than the sequential block accumulation it
    replaces, and the additional tree levels add kernel launch barriers.
\end{itemize}

\subsection{GPU Results}

\begin{table}[t]
\centering
\caption{NVIDIA B200 GPU execution time (ms), $Q = 256$, $k = n$.}
\label{tab:gpu}
\begin{tabular}{@{}rrr@{}}
\toprule
$n$ & Time (ms) & vs.\ CPU serial \\
\midrule
8{,}192    & [UPDATE] & [UPDATE]$\times$ \\
65{,}536   & [UPDATE] & [UPDATE]$\times$ \\
262{,}144  & [UPDATE] & [UPDATE]$\times$ \\
524{,}288  & [UPDATE] & [UPDATE]$\times$ \\
1{,}048{,}576 & [UPDATE] & [UPDATE]$\times$ \\
1{,}572{,}864 & 982 & [UPDATE]$\times$ \\
\bottomrule
\end{tabular}
\end{table}

The GPU frontier is $n = 1{,}572{,}864$ players at $k = n$ (all players
paid) in 982\,ms with $Q = 256$---71\% beyond the pre-optimization
baseline of $n = 917{,}504$ in 764\,ms. The frontier is limited not by
compute or memory bandwidth but by \emph{tree-level cliffs}: each
additional tree level increases compute by ${\sim}10\%$, creating
discrete jumps in execution time. The cliff at $n \approx 1{,}573{,}000$
is the first point where the tree's total work exceeds the one-second
budget.

[AUTHOR: Consider adding a figure showing execution time vs.\ $n$ with
the cliff structure visible, from the contour/heatmap data.]

% ======================================================================
\section{Conclusion}\label{sec:conclusion}
% ======================================================================

I have presented a high-performance algorithm for computing ICM tournament
equities using generating-function quadrature and FFT-accelerated
subproduct trees. The algorithm achieves $\bigO(Qn\log^2 k)$ time,
double-precision accuracy via $Q = 256$ quadrature points, and
automatically dispatches between a linear engine (optimal for small $k$)
and a hybrid block-tree engine (optimal for large $k$) using a roofline
cost model calibrated to the target hardware.

Three themes emerge from the implementation work. First,
\emph{the adjoint is the algorithm}: the propagation phase is precisely
the reverse-mode derivative of the build phase, and recognizing this
yields both a clean correctness proof and complexity parity for free.
Second, \emph{memory hierarchy dominates}: the roofline dispatch model,
the L2-aware checkpointing (whose omission in the parallel path caused
a 13.5$\times$ regression), and the GPU's FFT-stride layout all
demonstrate that data movement---not arithmetic---is the binding
constraint. Third, \emph{calibration beats modeling}: per-size FFT
timing, per-level cache bandwidths, and dual-library dispatch all exploit
the fact that measuring is cheaper and more accurate than predicting.

The GPU implementation on the NVIDIA B200 computes equities for
1,572,864 players in under one second, enabling ICM analysis at
field sizes that were previously intractable. Future work includes
extending the algorithm to non-uniform chip stack distributions with
adaptive quadrature, and exploring real-time tournament decision support
where subset equity queries can exploit the active-player masking already
present in my implementation.

\bibliographystyle{plainnat}
\bibliography{references}

\end{document}

%% ======================================================================
%% REVISION NOTES
%% ======================================================================
%%
%% Revision by: Claude (automated revision, 2026-03-29)
%% Based on: Author's revision instructions
%%
%% STRUCTURAL CHANGES:
%% 1. Added Section 2 (Related Work) covering: Malmuth-Harville model,
%%    bitmask DP, Monte Carlo ICM, polynomial subproduct trees (von zur
%%    Gathen & Gerhard), FFT libraries (FFTW, AOCL-FFTW, MKL), roofline
%%    model (Williams, Waterman & Patterson), automatic tuning (ATLAS,
%%    SPIRAL), GPU FFT (cuFFTDx, CUDA Graphs).
%%
%% 2. GPU section (now Section 6) completely rewritten: changed from H200
%%    roadmap to B200 results. Covers FFT-stride layout, cuFFTDx fused
%%    R2C/C2R, arena allocator, CUDA Graphs, Q-batch pipelining. Documents
%%    rejected optimizations (cuFFT callbacks: 2.3x slower; B=1 tree:
%%    20-56% slower). Frontier: n=1,572,864 at k=n in 982ms.
%%
%% 3. Engine dispatch (Section 4.6) rewritten as roofline cost model.
%%    Old Algorithm 5 (2nk * tau_poly * c_ckpt) replaced with
%%    bytes/blended_bw model using harmonic-mean bandwidth interpolation
%%    across L2->L3->DRAM. Validates at 42/44 correct decisions.
%%
%% 4. Added accuracy subsection (Section 5.3) comparing erfc-trap vs
%%    tanh-sinh quadrature across four distributions. erfc-trap converges
%%    to ~1e-12 by Q=256 uniformly; tanh-sinh gets lower floor on uniform
%%    but fails on geometric/adversarial.
%%
%% 5. Contributions list in intro updated: added roofline cost model and
%%    GPU implementation (B200, not H200).
%%
%% 6. Abstract shortened from ~200 words to ~150 words. Includes B200
%%    frontier result.
%%
%% CONTENT CHANGES:
%% 7. Zen 4 FFT backend changed from "FFTW+MKL" to "AOCL-FFTW+MKL".
%%    AOCL-FFTW (636 AVX-512 codelets vs 1 in upstream) delivers 1.2-2x
%%    at sizes 1K-32K. AOCL wins at 637/749 sizes, MKL at 112/749.
%%
%% 8. Parallel checkpointing fix documented: per-thread full-n allocation
%%    bug -> L2-aware ckpt in OpenMP path -> 13.5x speedup. Discussed in
%%    both the linear engine description and experimental results.
%%
%% 9. Removed detailed vDSP scaling fixup paragraph, AMX tile extraction
%%    details, and MKL dlopen mechanics. These are mentioned briefly in
%%    Related Work and platform descriptions but internals are not
%%    explained.
%%
%% 10. Fixed narrative arc: math insight (Sections 3-4) -> algorithmic
%%     challenge (Sections 4-5) -> systems challenge (Sections 5-6) ->
%%     results (Section 5 experiments + Section 6 GPU).
%%
%% 11. Added natbib package, \bibliography{references}, created
%%     references.bib with all citations.
%%
%% PLACEHOLDERS REMAINING:
%% - [UPDATE: re-measure ...] on Tables 1 and 2 (need final build numbers)
%% - [UPDATE] on several GPU table entries (need full GPU benchmark grid)
%% - [UPDATE] on M3 Max bandwidth constants in Table 3
%% - [AUTHOR: cite concrete MC-ICM paper] in Related Work
%% - [AUTHOR: fill in B200 SM count] in Experimental Setup
%% - [AUTHOR: Consider adding cliff figure] in GPU Results
%% - [AUTHOR: verify Knuth reference] in references.bib
%% - [AUTHOR: replace Chen 2010 placeholder] in references.bib
%%
%% STYLE:
%% - First person singular ("I") throughout
%% - Academic but direct, targeting SC/PPoPP venues
%% - Every claim backed by proof, citation, or [AUTHOR/UPDATE] marker
%% - No invented benchmark numbers
